% CLAUDE REPORT::BRIDGE-AUDIT-0001::2026-11-03
\documentclass{article}
\usepackage{amsmath,amssymb}
\begin{document}
\section*{Audit of the worst-case derivation bridge (BRIDGE-AUDIT-0001)}
Status: P versus NP remains OPEN. Nothing below is a new separation.

\paragraph{Corpus reviewed.} AGENTS.md, governance notes, the recent LSPV-1..4 reports on min-cost/min-DAG certificates (336--348), world-ledger pointer. No prior report states the two reductions below as a single audit.

\paragraph{Target.} $\overline W_R(N)=\max_{F\in\mathrm{UNSAT},|F|\le N}\min_{\pi:F\vdash_R\bot}\mathrm{cost}_R(\pi)$; bridge B: $\mathrm{SAT}\in P\Rightarrow \overline W_R\le \mathrm{poly}$.

\paragraph{Fact 1 (PROVED/EXACT).} If $\mathrm{cost}_R$ is proof size (or any measure polynomially related to the size of a verifiable certificate) and $R$ is a Cook--Reckhow system, then ``$\overline W_R$ superpolynomial'' is a statement about $R$ alone; free choice of the optimal $\pi$ is exactly nondeterminism. Superpolynomial $\overline W_R$ for \emph{every} such $R$ is equivalent to $\mathrm{NP}\neq\mathrm{coNP}$, strictly stronger than $P\ne NP$. For a single fixed $R$ it is the Cook--Reckhow program (open for Frege/EF).

\paragraph{Fact 2 (PROVED/EXACT, standard).} Bridge B for a fixed $R$ and proof-size cost fails for weak $R$ unconditionally in the relevant sense: resolution has exponential lower bounds on PHP/Tseitin while those families are in P. So $R$ must be strong enough that it p-simulates every polynomial-time SAT decider's trace.

\paragraph{Fact 3 (KNOWN THEOREM + argument).} A fixed uniform $R$ with B for \emph{all} polynomial algorithms must verify soundness of an unknown algorithm; Levin-style universal search certifies SAT witnesses but not UNSAT. Such an $R$ would, under $P=NP$, be a p-optimal-like system on the relevant fragment; existence of optimal/p-optimal systems is open (Krajicek--Pudlak 1989: implied by NE=coNE / E=NE) and Egidy (arXiv 2602.02294, Feb 2026) shows it is oracle-independent of PH structure. Hence B for a fixed $R$ cannot be derived by relativizing methods; it is itself an open optimality-type hypothesis.

\paragraph{Consequence.} The proposed route has two open arrows: (a) B (optimality-type, CANDIDATE), (b) superpolynomial $\overline W_R$ (at least as hard as a Cook--Reckhow lower bound for $R$). If cost includes simultaneous time and space of \emph{execution} rather than size, (a) is plausible via traces, but (b) then needs a genuine time--space lower bound for SAT, i.e. the ARC/FLMV spine already in the corpus. Recommended weakest next target: restrict $R$ to the \emph{trace system} $R_A$ of a fixed algorithm family and ask for $\overline W_{R_A}$ lower bounds only in a restricted-space model, linking to the TPORT/L13 barriers.

\paragraph{Literature clue.} Egidy 2026 (SPARSE-relativization) on (p-)optimal proof systems: relevant to what bridge B can and cannot rely on.

\paragraph{Stale/overstated claims.} None found in AGENTS.md; the directive itself correctly says no lower bound is established.
\end{document}
% END_OF_VERIFIED_REPORT::CLAUDE::BRIDGE-AUDIT-0001
